\chapter{Gonzales: A Production Path Tracer}
\label{chap:gonzales}

\section{Overview}

\gonzales{} is a physically based path tracer written from scratch in Swift~\cite{Wendleder2024Gonzales}.
It is, to our knowledge, the only single-author renderer capable of
processing Disney's \emph{Moana} island scene~\cite{Disney2018Moana}
in its entirety -- 130\,GB of geometry, textures, and shading data.

\begin{figure}[h]
  \centering
  \todo{Insert rendered Moana image here.}
  \caption{The Disney Moana island scene rendered by \gonzales{}.}
  \label{fig:moana}
\end{figure}

\section{System Architecture}

\subsection{Scene Loading and Parsing}

\todo{PBRT scene format, Ptex textures, instancing.}

\subsection{Acceleration Structure}

\gonzales{} uses Intel Embree~\cite{Wald2014Embree} for BVH construction
and ray-scene intersection.
Integration of Embree reduced render time of the Moana scene from
26~hours to 78~minutes.

\subsection{Material System}

\todo{Disney BRDF, multiple importance sampling, Veach 1997.}

\subsection{Integrator}

\todo{Path tracing integrator, Russian roulette, NEE.}

\section{The Moana Benchmark}

\todo{Scene statistics, validation against reference images, timing.}

\section{Implementation Notes}

\gonzales{} is implemented in approximately \todo{N} lines of Swift.
The choice of Swift over C++ provides memory safety, value-type semantics
for geometric primitives, and protocol-based polymorphism for materials
and integrators -- without sacrificing performance on hot paths.

\section{Summary}

\gonzales{} demonstrates that a single engineer can build a production-quality
renderer comparable in capability to team-built systems.
It serves as the software reference implementation for the integrated
\gonzales{}-on-\borg{} system described in \autoref{chap:gonzales_on_borg}.
